\documentclass[a4paper,12pt]{article}
\usepackage[russian]{babel}
\usepackage[utf8]{inputenc}
\usepackage[T2A]{fontenc}
\usepackage{geometry}
\usepackage{graphicx}
\usepackage{setspace}
\usepackage{anyfontsize}
\usepackage{float}
\usepackage{caption}
\usepackage{parskip}
\usepackage{enumitem}
\usepackage{hyperref}
\geometry{
a4paper,
left=20mm,
right=20mm,
top=15mm,
bottom=15mm
}
\begin{document}
\thispagestyle{empty}
\begin{titlepage}
  \begin{center}
    Министерство образования Республики Беларусь\\[1em]
    Учреждение образования\\
    БЕЛОРУССКИЙ ГОСУДАРСТВЕННЫЙ УНИВЕРСИТЕТ \\
    ИНФОРМАТИКИ И РАДИОЭЛЕКТРОНИКИ\\[1em]

    \begin{minipage}{\textwidth}
      \begin{flushleft}
        \begin{tabular}{ l l }
          Факультет & Информационных технологий и управления\\
          Кафедра   & Интеллектуальных информационных технологий
        \end{tabular}
      \end{flushleft}
    \end{minipage}\\[1em]

    \vspace{5em}


    \textbf{РАСЧЕТНАЯ РАБОТА}\\
    {по дисциплине <<Представление и обработка данных в интеллектуальных системах>>}\\
    на тему\\
    \textbf{\large Найти число вершинной связности для неориентированного графа.}\\[1em]

    \vspace{10em}
    
    \begin{tabular}{ p{0.65\textwidth}p{0.25\textwidth} }
      Выполнил:& П.\,Д.~Туровец \\[1em]
      Студент группы& \\
      421701 & \\
      Проверил: & Н.\,В.~Малиновская \\
     
    \end{tabular}
    
    \vfill
    {\normalsize Минск 2025}
  \end{center}
\end{titlepage}
\section{ВВЕДЕНИЕ}

\textbf{Цель: }Получить навыки формализации и обработки информации с
использованием семантических сетей

\textbf{Задача: }Найти число вершинной связности для неориентированного графа.

\section{СПИСОК ПОНЯТИЙ}
1. \textbf{\textit{Граф}} — математическая абстракция реальной системы любой природы, объекты которой обладают парными связями.
\begin{figure}[H]
\centering
\includegraphics[width=0.7\linewidth]{Граф.png}
\caption{Пример графа}
\label{fig:graph}
\end{figure}
2. \textbf{\textit{Неориентированный граф}} — это математическая структура, представляющая собой набор вершин и рёбер, где рёбра не имеют направления.
\begin{figure}[H]
\centering
\includegraphics[width=0.7\linewidth]{NotorientedGraph.png}
\caption{Пример неориентированного графа}
\end{figure}
\section{ТЕСТОВЫЕ ПРИМЕРЫ}

\subsection{Тест 1}
\begin{figure}[H]
\centering
\includegraphics[width=1.1\linewidth]{ex11.png}
\caption{Входные данные первого теста}
\includegraphics[width=0.8\linewidth]{ex12.png}
\caption{Выходные данные первого теста}
\end{figure}
\subsection{Тест 2}
\begin{figure}[H]
\centering
\includegraphics[width=1\linewidth]{ex21.png}
\caption{Входные данные второго теста}
\includegraphics[width=0.9\linewidth]{ex22.png}
\caption{Выходные данные второго теста}
\end{figure}
\subsection{Тест 3}
\begin{figure}[H]
\centering
\includegraphics[width=1.1\linewidth]{ex31.png}
\caption{Входные данные третьего теста}
\includegraphics[width=0.8\linewidth]{ex32.png}
\caption{Выходные данные третьего теста}
\end{figure}
\subsection{Тест 4}
\begin{figure}[H]
\centering
\includegraphics[width=1.1\linewidth]{ex41.png}
\caption{Входные данные четвёртого теста}
\includegraphics[width=0.8\linewidth]{ex42.png}
\caption{Выходные данные четвёртого теста}
\end{figure}
\begin{figure}[H]
\centering
\includegraphics[width=1\linewidth]{Формализация.png}
\caption{Формализация понятия "неориентированный граф"}
\end{figure}
\section{ПРИМЕР РАБОТЫ АЛГОРИТМА В СЕМАНТИЧЕСКОЙ ПАМЯТИ}
\begin{enumerate}[leftmargin=*,itemsep=2pt,parsep=2pt]
\item Шаг 1

Задаем неориентированный граф
\begin{figure}[H]
\centering
\includegraphics[width=0.8\linewidth]{pic1.png}
\caption{Неориентированный граф}
\end{figure}
\item Шаг 2

Перебираются все упорядоченные пары вершин (source, target) такие, что source < target. Рассмотрим первую пару: (A, B). Начинаем поиск пути между вершинами через BFS. Берём вершину A и заносим её в список visited.
\begin{figure}[H]
\centering
\includegraphics[width=0.8\linewidth]{pic2.png}
\caption{Начало обхода пары (А, B) с вершины А, занесение A в visited}
\end{figure}
\item Шаг 3

Проверяем вершину B: ребро A-B cуществует и B не посещена. Проверим и записываем новое значение числа вершнииой связности
\begin{figure}[H]
\centering
\includegraphics[width=0.8\linewidth]{pic3.png}
\caption{BFS в паре (A, B)}
\end{figure}
\item Шаг 4

Берём следующие пары. Остальные пары будут пропущены в формализации. Рассмотрим сразу пару (C, D), т.к. она даст полное понимание алгоритма. Начинаем поиск пути между вершинами через BFS. Берём вершину C и заносим её в список visited.   
\begin{figure}[H]
\centering
\includegraphics[width=0.8\linewidth]{pic4.png}
\caption{Занесение С в очищеный список visited}
\end{figure}
\item Шаг 5

Проверяем вершину A, она ещё не посещена и имеет ребро с C. Заносим её в visited.
\begin{figure}[H]
\centering

\includegraphics[width=0.8\linewidth]{pic5.png}
\caption{Проверка А}
\end{figure}
\item Шаг 6

Вершина B не имеет ребра с C, пропускае её. Идём к вершине C, пропускаем вершину, являющююся исходной. С вершиной D так же нет ребра, пропускаем и её. Все вершины проверены и мы переходим к следующей итерации: следующая вершина в очереди - A. Начинаем обходить с неё
\begin{figure}[H]
\centering
\includegraphics[width=0.8\linewidth]{pic2.png}
\caption{Начало новой итерации с А}
\end{figure}
\item Шаг 7

Проверяем вершину A, пропускаем связь "сам с собой". Переходим к вершине B: ребро с A имеется и B ещё не посещена. Заносим в visited и кладём в очередь.
\begin{figure}[H]
\centering
\includegraphics[width=0.8\linewidth]{pic6.png}
\caption{Проверка А}
\end{figure}
\item Шаг 8

Проверяем вершину C, уже посещяли, значит пропускаем. Смотрим на вершину D: ребро A-D существует, D ещё не посещена и является искомой. Проверим результат, так-же равен 1.
\begin{figure}[H]
\centering
\includegraphics[width=0.8\linewidth]{pic7.png}
\caption{Проверка D}
\end{figure}
\end{enumerate}

\section{ЗАКЛЮЧЕНИЕ}
В результате выполнения данной расчётной работы был формализован алгоритм построения каркаса графа, были изучены:
\begin{itemize}[leftmargin=*,itemsep=2pt,parsep=2pt]
\item Основы теории графов
\item Способы представления графов
\item Базовые алгоритмы для работы с графами
\item Основы SC-кода и SC-алфавита
\item Базовые знания о SCg-коде
\item Правила формирования идентификаторов sc-элементов
\end{itemize}
\section{ЛИТЕРАТУРА}
\par [1] Емеличев В. А., Мельников О. И., Сарванов В. И., Тышкевич Р. И. Лекции по теории
графов. М.: Наука, 1990. 384с. (Изд.2, испр. М.: УРСС, 2009. 392 с.)
\vspace{1em}
\par [2] Оре О. Теория графов. – 2-е изд.. – М.: Наука, 1980. – С. 336.
\vspace{1em}
\par [3] https://youtu.be/vMFwwgtEZSI?si=7rIjlNZu8TH0lCU
\end{document}
